% Small introduction, just to bring in the whitepaper
Designing a blockchain requires the careful consideration of numerous variables, each with its own unique, but impactful, influence on the final outcome of the network.
These considerations, in turn, are ossified into design decisions that set the trajectory of how the blockchain is governed, and how it will operate in the longer term.
This already difficult task of blockchain design, which requires a great appreciation for various disciplines of technology, economics, finance, philosophy and psychology, is now being swayed by the changing environment around us, coupled with increased usage and more complex interactions.

The nature of blockchain use is also growing, starting off as a niche, underground peer-to-peer digital currency, to now being embedded in global financial institutions and becoming an increasingly strong narrative for various interconnected applications of varying complexity.
This puts a degree of importance on designing blockchains, as both design and research should look beyond the current cycles with an attempt to visualise the potential that lies ahead.

% Then we can lead into the what we're bringing --> UT.
These new designs and ongoing research constantly push the boundaries of what is possible within blockchains, however, there are slight subtleties that inhibit their \emph{performance}, whether insufficient decentralization, low throughput, high transaction latency, weak protocol-level security assumptions, or, the lack of fostering native and intuitive interactions between communities and other blockchain networks.

One of the major challenges that has been identified to be a cause of such inhibitions is \emph{Buterin's Trilemma} (aka. the \emph{Scalability Trilemma})\footnote{%
  I think \emph{Scalability Trilemma} is a bad name because \emph{scalability} is one of the three conflicting qualities; you could just as well call it the \emph{Decentralization Trilemma}.
  Arguably, as we will see, the most appropriate of the 3 would actually be the \emph{Security Trilemma}.
  The term \emph{Blockchain Trilemma} has been used recently, which I think is worse.
  Apparently, the term was coined by Vitalik Buterin, so I prefer the name \emph{Buterin's Trilemma}.
},
which states that a blockchain can have at most two of the three following properties: (i)~security, (ii)~scalability and (iii)~decentralization (discussed below in \cref{subsec:trilemma}).

These properties, and the notion of only being able to ``\emph{choose two of three}'', quickly highlight potential shortcomings in various designs.
For example, Bitcoin and the original Proof-of-Work Ethereum (pre-merge) are chains that preference \emph{decentralization} and \emph{security} while foregoing scalability manifested into long transaction wait times, and higher fees due to the network not having enough capacity to process all requests.
On the other hand, blockchains that focus on \emph{security} and \emph{scalability} over decentralization often display exceptional transaction throughput, but are vulnerable to node failures causing network halts, or, fall to shortcomings of centralized services, such as censorship -- examples include Solana, BNBChain, Ripple, or enterprise blockchains.
Finally, chains that prefer \emph{scalability} and \emph{decentralization} over security may exhibit high performance, but may reduce parameters in their security choices for the sake of inclusivity, for example the original IOTA chain or chains that target IOT/embedded devices.

Amaroo's underlying consensus mechanism, \emph{Ultra Terminum} (UT), is a cooperative cross-chain consensus strategy that addresses \emph{Buterin's Trilemma}, and promotes cohesiveness amongst chains and communities.

%% Commenting for brevity of TODO
%% \mk{big claim, might have to qualify this a bit more about what 'existing consensus methods' means.}

Compared with existing consensus methods, UT provides \emph{equal or better} security properties
%% than \textbf{all} existing consensus methods
(including Bitcoin's PoW method and variants, and \emph{all} PoS variants).
This is because UT leverages existing consensus methods in combination and, by way of construction, UT can only \emph{add} security to these methods.
In a way, $\UT{}$ provides ways to generalize the contributions between chains, fostering collaboration to share security amongst heterogeneous chains.

\UT{} is capable of supporting millions of transactions per second with similar chain parameters (like block size) to those of traditional blockchains (like Bitcoin or Eth1\footnote{
    The term \emph{Eth1} refers to the original Proof-of-Work Ethereum (pre-merge), and \emph{Eth2} refers to the Sharded Proof-of-Stake Beacon Chain.
    These terms were originally adopted to describe the different designs on the \href{https://web.archive.org/web/20241128055038/https://ethereum.org/en/roadmap/}{Ethereum Roadmap}.
    }).
